\section{Extended Role Catalogue}
\label{app:roles}

The baseline five-role orthonormal basis in \cref{tab:roles} captures the
minimal functional partition observed across situationships. In larger social
portfolios ($n \ge 5$), finer specialization emerges. \cref{tab:extended-roles}
documents a ten-role extended basis alongside typical messaging indicators
and rank-confusion risk.

\begin{table}[htbp]
\centering
\footnotesize
\caption{Extended ten-role taxonomy for multi-partner relational portfolios.}
\label{tab:extended-roles}
\begin{tabularx}{\textwidth}{@{}lllL@{}}
\toprule
Symbol & Role name & Primary cadence & Characteristic observables and indicators \\
\midrule
$r_1$ & Public primary & Daily & Appears in public photographs; mutual family and friend recognition \\
$r_2$ & Nocturnal emotional & Late evening & Long voice notes after 11~p.m.; conflict debriefing \\
$r_3$ & Material/financial & Event-driven & Dinner invitations, errands, rides, practical favors \\
$r_4$ & Weekend leisure & Fri--Sun & Concerts, parties, spontaneous social outings \\
$r_5$ & Low-cadence reserve & Biweekly & Brief check-ins maintaining warm reactivation capacity \\
$r_6$ & Intellectual sparring & Daytime & Article sharing, career advice, academic collaboration \\
$r_7$ & Validation mirror & Variable & Frequent compliments, social-media story reactions \\
$r_8$ & Physical intimacy & Irregular & High-privacy meetings without public co-presence \\
$r_9$ & Travel/event escort & Seasonal & Dedicated companion for weddings, trips, or milestones \\
$r_{10}$ & Transitional buffer & Post-collapse & Rapidly elevated stand-in absorbing rebound attention \\
\bottomrule
\end{tabularx}
\end{table}

\begin{remark}[Basis reduction]
When modeling empirical logs with fewer than $m = 10$ distinct behavioral
clusters, roles $r_6$ through $r_{10}$ project cleanly onto linear combinations
of the canonical five-role subspace $\{r_1, \dots, r_5\}$ (for instance, $r_6$
partially aligns with $r_2$, while $r_9$ mixes $r_1$ and $r_4$).
\end{remark}
